\chapter{Introduction}
\label{chap:intro}

\section{Background and motivation}
\label{sec:motivation}

Portfolio optimisation allocates capital to different
assets in order to maximise returns for a given level of 
risk. Since Markowitz's seminal mean-variance model
\cite{markowitz1952portfolio}, the covariance
(or equivalently, correlation) matrix of returns has
been the main statistical object used for optimisation.
Hierarchical Risk Parity (HRP) is a recent
method \cite{lopezdeprado2016hrp} that extends the 
mean-variance literature and aims to address
the matrix instability problems that affected earlier 
models. HRP first clusters assets based on their 
correlational distance, producing a dendrogram (binary
tree), then recursively bisects down that tree and 
allocates capital to clusters inversely proportional
to their estimated variances.

However, correlation is an undirected metric: 
it records which assets co-move, but not which assets
move others \cite{lopezdeprado2023causal}. Crucially,
it lacks causal structure, so is prone to spurious
correlation \cite{simon1954spurious}.
Causal discovery algorithms provide a directed alternative
to the correlation matrix. Score-based structure learning algorithms
(such as NOTEARS \cite{zheng2018notears} and its temporal extension
DYNOTEARS \cite{pamfil2020dynotears}) can fit directed acyclic graphs
(DAGs) that can represent which assets are linked and the direction
of influence between them. Some papers have applied these DAGs to
equity markets \cite{howard2025causal,tang2024trading}, 
but do not use them for weighting or allocation within 
hierarchical portfolio allocators.

However, current hierarchical methods like HRP are not
able to represent directed objects like DAGs,
because HRP's clustering methodology requires symmetric matrices.
But a DAG can be separated into two information sources:
its skeleton, an undirected structure that displays which asset
pairs are causally linked; and its edge orientations, which indicate
the direction of each causal link. The skeleton can be directly
provided as an input to hierarchical allocators by symmetrising
the adjacency matrix representation of the graph. The edge
orientations cannot, as symmetrisation discards directionality,
so incorporating their information requires novel allocators.
In this thesis, I create these novel allocators (the
``orientation-aware'' allocator family) and create an experimental
design where each allocator receives the same DAG, 
so any performance difference is attributable to the information
(i.e. skeleton or skeleton + edge orientations) the allocator 
extracts from the DAG.

This research is informative for a variety of reasons.
First, running a causal discovery algorithm is computationally 
intensive, but once the DAG has been produced, 
using a skeleton in an existing hierarchical allocator is
trivial. Using the edge orientations' information, however,
requires the novel allocators developed in this research.
By decomposing a DAG into the two components,
and examining the marginal performance contributions of both, 
it informs practitioners as to whether the benefits of their
chosen portfolio optimisation approach are worth the costs.
Second, where previous research just constructed causal
graphs of assets \cite{howard2025causal}, this research
contributes to the causal finance literature by actually
setting portfolio weights based on the asset graphs,
and reporting the decomposed contribution of both aspects of
the graphs.
Third, the project demonstrates how small effects can be
presented comprehensively and honestly. Much
of the quantitative finance literature has a notable
problem with overfitting to backtests
\cite{bailey2014deflated,lopezdeprado2023causal}: when testing
a multitude of strategies, some will perform well by statistical
chance alone, though these chance outperformances are generally
very small in settings with long samples and highly correlated
strategies. I therefore distinguish between chance and real 
effects by reporting a variety of statistics, and 
accounting for the number of strategies searched over when 
reporting statistics.

\section{The research question}
\label{sec:question}

The thesis addresses a central research question:

\begin{quote}
When the
correlation matrix in hierarchical portfolio allocation is
replaced by a causal graph, is there a performance
improvement, and how much of any improvement is attributable
to the graph's skeleton versus its edge orientations?
\end{quote}

\noindent I briefly define the terms in the question here, and
expand on them in more detail later on in the report.
First, hierarchical portfolio allocation specifically refers to
HRP, where capital is invested fully and long-only. I invest over
a fixed 99-asset subset of the S\&P~500, with the portfolio
rebalanced monthly between 2007 and 2024. The asset ``universe''
is discussed in further detail in Chapter~\ref{chap:discovery},
the allocation protocol in Chapter~\ref{chap:evaluation}.
Second, the correlation matrix is the matrix originally used in HRP,
the lookback window's sample correlation
\cite{lopezdeprado2016hrp}. This feeds into the baseline HRP
control (see Chapter~\ref{chap:sensitivity}).
Third, a discovered causal graph is the output of one of the causal
discovery algorithms (DYNOTEARS, VARLiNGAM, with ridge-regularised
Granger as a robustness check). It is a contemporaneous (no lags)
asset-to-asset graph that's re-fitted at each rebalance
at 4 different estimation windows 
(189, 252, 378 and 504 trading days; roughly 9 months to 2 years).
See Chapter~\ref{chap:discovery} for more details.
Fourth, performance improvement is defined as the out-of-sample 
annualised Sharpe ratio of daily portfolio returns, net of 
investment cost. I also use several other metrics that account 
for multiplicity (see Chapter~\ref{chap:evaluation}), including the 
Deflated Sharpe ratio, Hansen's SPA, and the Model Confidence Set
\cite{bailey2014deflated,hansen2005spa,hansen2011mcs}.
Finally, I compare the relative performance gain attributable to
the graph's skeleton versus its edge orientations using a 
``fixed-graph ablation'', where an identical graph $M$ feeds
a variety of allocators that differ only in the information they
extract from the graph. The skeleton-only allocator, \skelsamp{},
symmetrises $M$ before clustering, while three 
orientation-aware allocators consume the directed matrix itself
(either through the allocation step, the ordering step, or both).
See Chapter~\ref{chap:sensitivity} for further details.

The research question contains two requirements, both of which I test:
\begin{enumerate}[itemsep=1pt, topsep=3pt]
    \item The causal graph must improve on the correlation matrix
    for portfolio allocation. This is the first part of the question,
    which then leads to the second part of the question regarding the
    relative contribution from the skeleton and the edge orientations
    of the graph. In any case, even if this requirement fails, the
    decomposition of the relative performance of the skeleton and
    the edge orientations is informative. We test the performance
    by providing the HRP allocator with the symmetrised causal graph
    (this allocator is denoted \skelsamp{} in the thesis), and 
    comparing it to the standard, sample correlation-based HRP method,
    in Section~\ref{sec:r1}.
    \item Allocators must be able to use directional information
    to take advantage of edge orientations. Symmetric clustering,
    which HRP uses, cannot, so I construct an orientation-aware
    family of allocators in Chapter~\ref{chap:sensitivity}.
    I also verify that the skeleton-only allocators and the
    HRP baseline are unable to take in orientation information 
    (see Chapter~\ref{chap:sensitivity}).
\end{enumerate}

\noindent Assuming both of the requirements are met, I am then able
to isolate the performance gain of the two components of the graph.
The skeleton's contribution over the HRP baseline measures the
value of knowing which assets are causally linked to each other. 
The edge orientations' contribution over the skeleton measures
the additional value of knowing the direction of the causal links
between assets, holding the graph fixed. The two marginal
contributions sum to the total performance gain of the causal
graph over the correlation-based allocator. 

A subsidiary question follows once the decomposition is measured: through
what channel does orientation pay, directional content read from the
arrows, or properties of the implied covariance that a graph-free
construction can reproduce? The distinction matters, because only the
first would vindicate a causal reading of the graph.
Chapter~\ref{chap:evaluation} answers it with two covariance controls,
Ledoit--Wolf
shrinkage and a de-factored covariance, and completes the symmetry with a
graph-free skeleton control that asks the same question of the skeleton
channel.

\section{Challenges and contributions}
\label{sec:challenges}

The work is organised as three technical challenges (C1--C3) and the
deliverables addressing them (A1--A3), one chapter each. The mapping onto the
research question
is direct: C1 supplies the graph, C2 supplies the allocators that decompose
its value, and C3 supplies the inference that adjudicates the decomposition.
Together they integrate three literatures that have not previously met at
the allocation step (causal structure learning, hierarchical portfolio
construction, and the backtest-rigour programme), and A1--A3 are the
project's objectives: Chapters~\ref{chap:discovery}--\ref{chap:evaluation}
deliver them, and Chapter~\ref{chap:evaluation} assesses them jointly
against the research question.

\begin{enumerate}[label=C\arabic*, itemsep=2pt, topsep=3pt]
  \item Discovering a credible asset--asset causal graph at every
    rebalance. The graph must be estimated over a ${\sim}130$-variable joint
    panel of assets and macro drivers, under an economically-motivated
    directional prior (drivers may cause assets; single stocks do not cause
    the 10-year Treasury yield), 215 times per configuration, cheaply enough
    to sweep, and reproducibly: a graph that changes between runs
    invalidates every downstream comparison.
  \item Feeding a directed graph to an allocator that fundamentally
    consumes symmetric matrices. Hierarchical clustering requires a
    symmetric distance; naive symmetrisation discards exactly the
    orientations in question. Exploiting them requires allocation machinery
    with no analogue in the hierarchical-allocation literature, and an
    experimental design
    in which orientation is the only treatment variable.
  \item Attributing small performance differences credibly. The
    candidate effects are $0.01$--$0.03$ Sharpe on heavy-tailed returns,
    produced by a study that evaluated \rsNtrials\ configurations, so
    the search itself, if unpriced, could manufacture the whole effect.
\end{enumerate}

\begin{enumerate}[label=A\arabic*, itemsep=2pt, topsep=3pt]
  \item A reproducible rolling pipeline for causal discovery
    (addresses C1; Chapter~\ref{chap:discovery}): DYNOTEARS with the directional prior
    enforced natively and empirically verified to matter; a content-keyed
    discovery cache that makes the full ablation grid tractable
    (a backtest re-runs in ${\approx}40$ seconds against cached fits
    instead of ${\approx}15$ hours); a data-determinism defect found
    and fixed; and a single extraction chokepoint guaranteeing that every
    allocator receives the identical graph.
  \item The orientation-aware allocator family and the fixed-graph
    ablation (addresses C2; Chapter~\ref{chap:sensitivity}), the core
    methodological novelty: a covariance implied by a structural equation
    model (SEM), which propagates shocks through all directed paths of the
    DAG, and
    causal-ordered bisection, which replaces the clustering
    dendrogram with the DAG's topological order. Each is deterministic and
    seed-free, and each is evaluated against matched skeleton-only and
    correlation controls on the identical graphs.
  \item A multiplicity-, distribution- and stability-aware evaluation
    (addresses C3; Chapter~\ref{chap:evaluation}): stationary-block-bootstrap
    contrasts for every pairwise step of the decomposition; and the
    PSR/DSR/Reality-Check/SPA/MCS battery (the DSR deflating against all
    \rsNtrials\ trials, the family-wise tests over their stated
    per-question families).
\end{enumerate}

\section{Quantified benefits}
\label{sec:benefits}

All benefits below are measured on the protocol of
Section~\ref{sec:question} (net of 5\,bps, monthly rebalances, 2007--2024)
with significance from the Politis--Romano stationary block bootstrap
\cite{politis1994stationary}, and are reported as measured, including
where the answer is small or null. Full detail is in
Chapter~\ref{chap:evaluation}.

\begin{itemize}[itemsep=2pt, topsep=3pt]
  \item The decomposition (the answer to the research question).
    The graph's total gain over the correlation control is stable across
    estimation windows ($+0.018$ to $+0.032$ net Sharpe for \skelsem{}), worth
    $+0.4$ to $+0.7$ percentage points of net compound annual growth rate
    (CAGR) per year, but its composition shifts with the window: the
    skeleton's contribution is hump-shaped and peaks near the conventional
    one-year window, while orientation's is positive at every window and
    largest at the extremes, though the between-window differences are
    about one standard error and the window grid was widened partway
    through. Statistical support is modest: one pairwise contrast in
    twelve is individually significant, and family-wise support exists
    only under a narrowed family fixed after the results were seen
    (Chapter~\ref{chap:evaluation} gives the full estimates, the
    anchor-robustness caveats and the multiplicity treatment).
  \item What orientation buys is robustness, not level. The
    orientation-aware allocators are strong at every window (\skelsem{}
    spans $0.395$--$0.415$ and \orientsem{} $0.399$--$0.419$ across the
    four estimation windows, against the skeleton-only allocator's
    $0.372$--$0.403$), and their edge over the skeleton is largest where
    the sample moments feeding the control are weakest (the nine-month and
    two-year windows), consistent with structure substituting for
    sample-moment quality. A direct control identifies the mechanism: the
    SEM covariance's removal of the common market factor, which a
    graph-free de-factored covariance reproduces at every window
    (Chapter~\ref{chap:evaluation}), so the robustness is real but is not
    evidence of directional content. A matching graph-free skeleton
    control shows the converse for the skeleton channel: its one-year
    gain does require the discovered graph. At the 18-month window
    \orientsem{} attains the highest Deflated
    Sharpe ratio of all \rsNtrials\ configurations
    ($\rsDtwosDSRWthree$). The effect is, however,
    method-dependent (absent under VARLiNGAM graphs at all four windows)
    and, descriptively, concentrates in calm regimes (the bottom quintile
    of the Cboe Volatility Index, VIX) rather than stress.
  \item Novelty relative to the closest published work. Howard et
    al.\ \cite{howard2025causal} build causal graphs on the S\&P~500 but
    never set weights from them. Setting hierarchical weights
    directly from a causal graph, decomposing skeleton from orientation on
    fixed graphs, and the orientation-aware allocator family itself are, to
    my knowledge, new. Because no published configuration sets hierarchical
    weights from a discovered graph, no external number exists for a direct
    comparison; the external benchmark is therefore the \HRP{} control
    itself, L\'opez de Prado's published textbook construction, evaluated
    on the identical protocol.
\end{itemize}

In summary: the causal graph's value to hierarchical allocation
is stable across estimation windows, but its source shifts with them,
and family-wise significance exists only under the post-hoc narrowed
family. The contribution is the allocator family
and the fixed-graph design that made this measurable;
Chapter~\ref{chap:conclusion} restates the full answer with its caveats.

\section{Report structure}
\label{sec:structure}

Chapter~\ref{chap:background} analyses the two works this study builds
on (HRP and the backtest-rigour agenda, and causal graphs in equity
markets) and locates the gap. Chapter~\ref{chap:discovery} (A1) presents the
discovery pipeline that produces the graphs. Chapter~\ref{chap:sensitivity}
(A2) develops the allocator family and the fixed-graph ablation that make the
research question
answerable. Chapter~\ref{chap:evaluation} (A3) evaluates: the two
requirements, the decomposition itself, and how much of it survives rigorous
inference. Chapter~\ref{chap:conclusion} concludes with limitations,
broader impact, and future work. Appendix~\ref{app:tables} collects the full
result tables,
operative hyperparameters and reproducibility notes; the report stands
without it.
